\documentclass[10pt]{article}
\input{../../notes-style.tex}

\usepackage{multicol}
\usepackage{array}

\title{Ch.~3 Cheat Sheet\\\large Data Structures, Big-O, Sorting}
\author{}
\date{}

\begin{document}
\maketitle
\vspace{-2em}

\begin{multicols}{2}

\subsection*{ADT vs.\ data structure}
\textbf{ADT} = interface (\emph{what}). \textbf{Data structure} =
implementation (\emph{how}). One ADT, many implementations.
\begin{tabular}{@{}ll@{}}
\textbf{ADT}           & \textbf{typical DS}\\
List                   & array, linked list\\
Stack                  & array / linked list\\
Queue                  & linked list / ring buf\\
Priority queue         & binary heap\\
Set                    & hash table / BST\\
Map / Dictionary       & hash table / BST\\
Graph                  & adj.\ list / matrix
\end{tabular}

\subsection*{The seven structures}
\begin{itemize}
    \item \textbf{Array / vector} -- contiguous; $O(1)$ indexing.
    \item \textbf{Linked list} -- nodes + pointers; $O(1)$ insert if you
    hold the prev.
    \item \textbf{Tree} -- hierarchical; recursion natural.
    \item \textbf{BST} -- ordered tree; $O(\log n)$ ops if balanced.
    \item \textbf{Heap} -- complete binary tree; $O(1)$ peek, $O(\log n)$
    push/pop.
    \item \textbf{Hash table} -- bucketed; $O(1)$ expected ops.
    \item \textbf{Graph} -- nodes + edges; adjacency list or matrix.
\end{itemize}

\subsection*{Searching}
\begin{tabular}{@{}lll@{}}
Linear search & any data     & $O(n)$\\
Binary search & sorted       & $O(\log n)$
\end{tabular}
Binary search midpoint: \texttt{lo + (hi - lo) / 2} (avoid overflow).

\subsection*{Big-O in one sentence}
$f(n) = O(g(n))$ iff $\exists c > 0, n_0$ such that $f(n) \leq c \cdot g(n)$
for all $n \geq n_0$. Drop constants, drop lower-order terms, keep the
dominant term.

\subsection*{Growth hierarchy}
\[
O(1) < O(\log n) < O(n) < O(n \log n) < O(n^2) < O(2^n) < O(n!)
\]
\textbf{Rule of thumb for loops:}
\begin{itemize}
    \item Single loop over $n$: $O(n)$.
    \item Nested over $n$: $O(n^2)$.
    \item Halving each iteration: $O(\log n)$.
    \item Divide-and-conquer ($T(n) = 2T(n/2) + O(n)$): $O(n \log n)$.
\end{itemize}

\subsection*{Constant-time operations}
\texttt{v[i]}, \texttt{v.back()}, \texttt{push\_back} (amortized),
\texttt{pop\_back}, \texttt{size()}, arithmetic, a bounded-iteration
loop. \emph{Not} constant: hashing a long string, a loop whose bound
depends on input.

\columnbreak

\subsection*{Sorting comparison}
\begin{tabular}{@{}l|ccc|cc@{}}
\textbf{Sort}        & \textbf{best}     & \textbf{avg}      & \textbf{worst}    & \textbf{stable} & \textbf{in-place}\\
\hline
Bubble               & $n$               & $n^2$             & $n^2$             & yes             & yes\\
Insertion            & $n$               & $n^2$             & $n^2$             & yes             & yes\\
Selection            & $n^2$             & $n^2$             & $n^2$             & no              & yes\\
Shell                & $n \log n$        & var.              & $n^{3/2}$         & no              & yes\\
Merge                & $n \log n$        & $n \log n$        & $n \log n$        & yes             & no\\
Quick                & $n \log n$        & $n \log n$        & $n^2$             & no              & yes\\
Heap                 & $n \log n$        & $n \log n$        & $n \log n$        & no              & yes\\
Radix (LSD)          & $nk$              & $nk$              & $nk$              & yes             & no
\end{tabular}

\subsection*{``Which sort do I pick?''}
\begin{itemize}
    \item \textbf{Small array ($n < 20$)} or \textbf{nearly sorted}:
    insertion sort ($O(n)$ best case, tiny constants).
    \item \textbf{General-purpose, RAM-bound}: quicksort (what
    \texttt{std::sort} uses, with heapsort fallback).
    \item \textbf{Stability required}: merge sort (\texttt{std::stable\_sort}).
    \item \textbf{External / too big for RAM}: merge sort (sequential
    I/O; the model for external mergesort).
    \item \textbf{Integer keys, bounded width}: radix sort (beats
    comparison-based $\Omega(n \log n)$ bound).
    \item \textbf{Tightest worst-case guarantee}: heap sort.
\end{itemize}

\subsection*{The $\Omega(n \log n)$ lower bound}
Any \emph{comparison-based} sort needs $\Omega(n \log n)$ comparisons in
the worst case. Decision-tree argument: $n!$ possible outputs, so any
binary tree that distinguishes them has depth $\geq \log_2(n!) = \Theta(n
\log n)$. Radix and counting sort beat this because they're
\emph{non-comparison}.

\subsection*{Stable vs.\ unstable}
A sort is \textbf{stable} if equal keys retain their original relative
order. Matters when sorting by one field after another (e.g., sort by
last name, then stable-sort by city).

\subsection*{In-place}
Uses $O(1)$ or $O(\log n)$ auxiliary memory. Merge sort's $O(n)$ buffer
makes it \emph{not} in-place; its recursion stack is $O(\log n)$.

\subsection*{Quicksort: watch the pivot}
Average $O(n \log n)$; worst $O(n^2)$ on already-sorted input with
first-element pivot. Fixes: median-of-three; randomize. \texttt{std::sort}
uses introsort: quicksort + fallback to heapsort after a depth limit.

\subsection*{Top gotchas}
\begin{itemize}
    \item Quoting ``$O(n \log n)$ always'' for quicksort (it's average;
    worst case is $O(n^2)$).
    \item Saying merge sort is in-place (it's not).
    \item Using \texttt{strcmp}-style comparators for numeric sort.
    \item Assuming \texttt{std::sort} is stable (use
    \texttt{std::stable\_sort}).
    \item Reasoning about ``fast'' without bounding $n$.
\end{itemize}

\end{multicols}
\end{document}
